\documentclass[11pt]{article}
\usepackage[margin=1in]{geometry}
\usepackage{amsmath, amssymb}

\title{Exercise 3.16: Adding a Constant to All Rewards (Episodic Tasks)}
\author{Sutton \& Barto, Section 3.5}
\date{}

\begin{document}
\maketitle

\section*{Problem}

Now consider adding a constant $c$ to all the rewards in an episodic task,
such as maze running. Would this have any effect, or would it leave the task
unchanged as in the continuing task above? Why or why not? Give an example.

\section*{Solution}

\subsection*{Episodic return with shifted rewards}

In an episodic task, the return from time $t$ is a \emph{finite} sum:
\[
    G_t = \sum_{k=0}^{T-t-1} \gamma^k R_{t+k+1}
\]

Adding $c$ to every reward:
\begin{align*}
    G'_t &= \sum_{k=0}^{T-t-1} \gamma^k (R_{t+k+1} + c)
         = G_t + c \sum_{k=0}^{T-t-1} \gamma^k
\end{align*}

\subsection*{Why this changes the task}

The added term $c \sum_{k=0}^{T-t-1} \gamma^k$ depends on $T - t$, the number
of steps remaining in the episode. Unlike the continuing case (where $\sum \gamma^k
= \frac{1}{1-\gamma}$ is a fixed constant), the episode length $T$ varies across
episodes and depends on the policy. Therefore:

\begin{itemize}
    \item The shift is \textbf{not} the same for all states or all policies.
    \item Adding $c$ \textbf{does} change the relative values of states and can
          change which policy is optimal.
\end{itemize}

\subsection*{Example: Maze running}

Consider a maze where rewards are $-1$ per step (encouraging the agent to
escape quickly). Adding $c = +2$ to all rewards makes each step yield $+1$.

\begin{itemize}
    \item \textbf{Before} ($R = -1$ per step): The agent wants to reach the
          exit as quickly as possible to minimize the total penalty.
    \item \textbf{After} ($R = +1$ per step): The agent is now \emph{rewarded}
          for each step. Longer episodes yield higher returns, so the agent is
          incentivized to \emph{avoid} the exit and wander as long as possible.
\end{itemize}

The optimal policy is completely reversed. This shows that, unlike in continuing
tasks, the absolute values of rewards (not just their relative differences)
matter in episodic tasks.

\subsection*{Why the continuing case was immune}

In the continuing case, every policy generates an infinite stream of future
rewards. The constant $c$ contributes $\frac{c}{1-\gamma}$ regardless of the
policy, so it cancels out when comparing policies. In the episodic case,
different policies produce episodes of different lengths, so the contribution
of $c$ varies and does not cancel.

\end{document}
